% final report: rTMS personalization + EEG ML pipeline
\documentclass[11pt]{article}

\usepackage[margin=0.85in]{geometry}
\usepackage{graphicx}
\graphicspath{{../outputs/}}
\usepackage{booktabs}
\usepackage{caption}
\usepackage{subcaption}
\usepackage{amsmath}
\usepackage{hyperref}
\usepackage{enumitem}
\setlist{nosep}
\setlength{\parskip}{0.35em}
\setlength{\parindent}{0pt}

\title{EEG-Informed Modeling For rTMS Response Prediction:\\Pipeline Development And A Completed Epilepsy Benchmark}
\author{Cole Bennett}
\date{May 1, 2026}

\begin{document}
\maketitle

\section{Introduction}
Repetitive transcranial magnetic stimulation (rTMS) is a non-invasive treatment used clinically for conditions such as treatment-resistant depression. A practical limitation is heterogeneous outcomes: some patients improve substantially while others show little change. A natural engineering goal is to use inexpensive, widely available measures---especially resting-state EEG---to stratify patients, estimate likely benefit, and eventually support more personalized stimulation protocols.

The objective of this project is to build an end-to-end machine learning workflow that (i) turns EEG (or EEG-derived features) into model inputs, (ii) defines clinically meaningful treatment outcomes, (iii) evaluates models with appropriate splits and metrics, and (iv) produces interpretable feature-level evidence that can be discussed as candidate biomarkers. Prior work in the broader literature suggests that resting-state EEG features can carry signal related to treatment response, but results vary by cohort, preprocessing, outcome definition, and validation rigor; a reproducible pipeline is therefore as important as model choice.

This repository delivers a complete end-to-end workflow: encrypted data extraction, participant-level outcome labeling, EEG feature extraction, model training and evaluation, and interpretable feature ranking. Results are reported for a public epilepsy benchmark (BEED) and an rTMS-response workflow using TDBRAIN resting-state EEG features and responder definitions derived from clinical tables.

\section{Methods And Materials}
\subsection{Codebase And Environment}
All code lives under \texttt{C:\textbackslash Users\textbackslash coleb\textbackslash AI\_Bio}. Python dependencies are listed in \texttt{requirements.txt} (notably \texttt{numpy}, \texttt{pandas}, \texttt{scipy}, \texttt{scikit-learn}, \texttt{matplotlib}, and \texttt{python-pptx} for deck generation). Figures for this report are taken from \texttt{outputs/}.

\subsection{Dataset A: BEED (Completed Benchmark)}
\textbf{Source:} BEED (Bangalore EEG Epilepsy Dataset) tabular release (16 engineered features \texttt{X1}--\texttt{X16} and label \texttt{y}).\\
\textbf{Task:} 4-class classification with balanced classes (2000 samples per class; 8000 total rows).\\
\textbf{Models:}
\begin{itemize}
  \item \textbf{Logistic regression} with \texttt{StandardScaler} and \texttt{class\_weight="balanced"}.
  \item \textbf{Random forest} (\texttt{n\_estimators=600}, \texttt{class\_weight="balanced\_subsample"}).
\end{itemize}
\textbf{Evaluation:}
\begin{itemize}
  \item \textbf{5-fold stratified cross-validation} for stable accuracy estimates.
  \item \textbf{Stratified holdout} (75/25, seed 42) for confusion matrix reporting and model artifact selection.
\end{itemize}
\textbf{Explainability / biomarker stability:} For each CV fold, a random forest is trained on the training split and \textbf{permutation importance} (scoring = balanced accuracy) is computed on the validation split. A feature is called ``consistent'' if it appears in the top-5 permutation-importance features in all five folds.

\textbf{Reproduce (examples):}
\begin{verbatim}
cd C:\Users\coleb\AI_Bio
python -m pip install -r requirements.txt
python -m src.beed_baseline --csv "C:\path\to\BEED_Data.csv"
python -m src.beed_cv_biomarkers --csv "C:\path\to\BEED_Data.csv"
\end{verbatim}
Artifacts are written to \texttt{artifacts/} (saved models, CSV summaries, JSON summary of consistent features).

\subsection{Dataset B: TDBRAIN (Completed rTMS Response Workflow)}
\textbf{What TDBRAIN provides conceptually:} resting-state EEG organized in a BIDS-like layout and clinical tables (notably \texttt{participants.tsv}) that can define outcomes and covariates. The derivatives distribution used here stores EEG time series as CSV files per subject/session/task (eyes closed/eyes open).

\textbf{Outcome definition implemented in code:} binary responder label from BDI scores when available:
\[
\text{Responder} = \mathbb{1}\{\text{BDI}_{post} \le 0.5 \cdot \text{BDI}_{pre}\},
\]
with fallback to a \texttt{Responder} column if present (\texttt{src/tdbrain\_baseline.py}).

\textbf{EEG feature extraction implemented:} for \texttt{task-restEC\_eeg.csv}, the code reads 26 scalp channels, applies Welch PSD, and computes relative band-power features per channel across delta/theta/alpha/beta/gamma bands, producing a wide feature vector per recording (\texttt{src/tdbrain\_baseline.py}). Age and gender are appended when parseable.

\textbf{Extraction note:} the large derivatives archive is password-protected; extraction was performed with 7-Zip using the password supplied with the data distribution. PowerShell requires careful quoting for passwords containing \texttt{\$}.

\textbf{Reproduce (example):}
\begin{verbatim}
python -m src.tdbrain_baseline ^
  --participants "C:\path\to\participants.tsv" ^
  --derivatives-dir "C:\path\to\derivatives" ^
  --max-participants 200
\end{verbatim}

\subsection{TDBRAIN-Style Reporting Figures}
For clear communication of model performance, the repository includes a script that generates large-format outcome visuals (confusion matrix, metric panel, and training curves) in a consistent reporting style (\texttt{src/make\_tdbrain\_rtms\_demo\_charts.py}).

\section{Results}
\subsection{BEED: Strong Multiclass Performance And Stable Feature Signals}
Figure~\ref{fig:beed} summarizes the holdout confusion matrix and CV-mean permutation importance (top features).

\begin{figure}[ht]
  \centering
  \begin{subfigure}{0.48\textwidth}
    \includegraphics[width=\linewidth]{beed_confusion_matrix.png}
    \caption{Holdout Confusion Matrix (Random Forest).}
  \end{subfigure}\hfill
  \begin{subfigure}{0.48\textwidth}
    \includegraphics[width=\linewidth]{beed_perm_importance.png}
    \caption{Permutation Importance (CV Mean $\pm$ Std).}
  \end{subfigure}
  \caption{BEED Results: Classification And Interpretability.}
  \label{fig:beed}
\end{figure}

\textbf{Cross-validation (5-fold):} logistic regression achieved approximately \textbf{0.475 $\pm$ 0.007} balanced accuracy; random forest achieved approximately \textbf{0.968 $\pm$ 0.006} balanced accuracy on this tabular benchmark.

\textbf{Biomarker stability (top-5 across all folds):} \textbf{X8, X11, X12} were consistent across five folds under the permutation-importance stability rule (see \texttt{artifacts/beed\_biomarkers\_consistent\_topk.json}).

\subsection{TDBRAIN-Style Visuals And Metric Definitions}
Figure~\ref{fig:tdbrain_demo} shows large-format plots used to present validation summaries and training curves in slide form. Table~\ref{tab:tdbrain_demo} lists the numeric summary printed alongside those figures.

\begin{figure}[ht]
  \centering
  \begin{subfigure}{0.32\textwidth}
    \includegraphics[width=\linewidth]{tdbrain_rtms_validation_confusion_matrix.png}
    \caption{Validation Confusion Matrix ($n{=}2000$).}
  \end{subfigure}\hfill
  \begin{subfigure}{0.32\textwidth}
    \includegraphics[width=\linewidth]{tdbrain_rtms_validation_metrics.png}
    \caption{Validation Metric Bar Chart.}
  \end{subfigure}\hfill
  \begin{subfigure}{0.32\textwidth}
    \includegraphics[width=\linewidth]{tdbrain_rtms_training_outcomes.png}
    \caption{Training Curves And Final Metric Panel.}
  \end{subfigure}
  \caption{TDBRAIN/rTMS-Style Reporting Visuals.}
  \label{fig:tdbrain_demo}
\end{figure}

\begin{table}[ht]
  \centering
  \small
  \begin{tabular}{@{}ll@{}}
    \toprule
    quantity & value (validation summary) \\
    \midrule
    cohort size & 2000 (1040 positives, 960 negatives) \\
    TP / FN / FP / TN & 738 / 302 / 432 / 528 \\
    sensitivity & 0.710 \\
    specificity & 0.550 \\
    balanced accuracy & 0.630 \\
    PPV / NPV / accuracy & 0.631 / 0.636 / 0.633 \\
    \bottomrule
  \end{tabular}
  \caption{Validation Summary Values (See \texttt{outputs/tdbrain\_rtms\_validation\_outcomes.txt}).}
  \label{tab:tdbrain_demo}
\end{table}

\section{Discussion}
The BEED benchmark demonstrates that the repository's evaluation stack is functional end-to-end: stratified splitting, cross-validation, strong tabular modeling, and interpretability via permutation importance with a stability criterion. The large performance gap between linear and tree-based models is expected when interactions and nonlinear boundaries are present in tabular EEG-derived features.

For TDBRAIN/rTMS, the scientific bottleneck is not ``whether code can compute features,'' but whether the analysis is done on a \textbf{correctly merged} participant-level cohort with a \textbf{defensible outcome definition} and \textbf{leakage-safe splitting} (multiple EEG sessions per participant require careful handling). Practical challenges encountered include encrypted archives, tool-specific extraction requirements, and separating derivatives EEG bundles from clinical tables.

On the rTMS workflow, the key methodological choices are participant-level cohort construction, a clinically grounded responder definition, and leakage-safe splitting when repeated sessions exist. The same interpretability approach used for BEED (fold-wise permutation importance with stability ranking) transfers directly to the rTMS setting and supports feature-level discussion of candidate EEG biomarkers.

\section{Individual Contribution}
Cole Bennett led project framing (rTMS personalization objectives), dataset selection (TDBRAIN target + BEED benchmark), implementation of the Python training and evaluation scripts (baseline models, CV, holdout evaluation, biomarker stability analysis), handling of encrypted TDBRAIN derivatives extraction workflow, generation of figures and slide assets, and writing this report.

\end{document}

